\section{Conclusion and Outlook}
\label{sec:conclusion}

This thesis set out to determine whether quantum machine learning offers a practical
advantage for the classification of astronomical spectra, and to answer that question
honestly rather than optimistically. The guiding principle throughout was fairness.
Every quantum model was compared against a classical counterpart of matched size,
trained under an identical pipeline, so that any measured difference could be
attributed to the quantum component rather than to incidental design choices.

We first established a strong classical foundation. A 14-million-parameter
convolutional network reached 84.5\% accuracy on the full 62-class SDSS problem and
reproduced the GaSNet-II benchmark at 92.5\% on a 13-class subset, confirming that the
classical baseline is competitive with the published state of the art and that the
coarse classification problem is essentially solved. This baseline also served as the
frozen feature extractor for the later experiments.

On top of this foundation we tested four different quantum approaches. A variational
quantum circuit on the hardest binary pair reached the same accuracy as its
parameter-matched classical mirror. Quantum and classical heads on frozen features both
reached about 96\% on a four-class task, with no parameter-efficiency benefit for the
quantum head. Quantum support vector machines showed no advantage in the few-shot
regime and were overtaken by a classical kernel as the training set grew. A
quanvolutional filter was statistically indistinguishable from a parameter-matched
classical filter under increasing noise, and the only significant difference favoured
the classical model. Across every paradigm the result was the same, namely parity at
best and no quantum advantage.

Just as important as the null result is how it was reached. In two separate experiments
an apparent quantum advantage appeared and then dissolved under scrutiny, once traced
to an under-built classical baseline and once to a dead-ReLU optimisation artefact that
a single change of activation removed. These episodes are the heart of the thesis's
contribution. The work delivers a disciplined fair-comparison methodology, with
parameter matching, shared pipelines, robust classical controls, multiple seeds and
statistical testing, and it demonstrates concretely that apparent quantum advantages
must be treated as hypotheses to be ruled out rather than as results. Because all
circuits were simulated without noise, the most favourable possible setting for quantum
hardware, the absence of an advantage is a conservative and robust conclusion. For the
spectral classification tasks studied here, and at the scales reachable on near-term
devices, quantum machine learning matched but did not beat a carefully built classical
model.

\subsection{Outlook}
\label{sec:conclusion_outlook}

Within this thesis the most obvious quantum-specific levers were already pulled. We
varied the circuit geometry across the full sweep and we compared angle encoding
against amplitude encoding, and neither the shape of the circuit nor the way the data
entered it produced an advantage over a matched classical model. The more promising
open questions therefore lie less in tuning these hybrid models and more in changing the
problem they are asked to solve.

The clearest of these is whether quantum methods help at all on data that is already
well suited to classical learning. Optical spectra are smooth, structured signals that
convolutional networks handle naturally, which leaves little room for a quantum model to
add anything. A fairer test of a genuine quantum advantage may require problems whose
structure is provably hard for classical learners, or data that is quantum in origin
rather than a classical measurement, where a quantum model is not forced to first
compress the input through a classical front end. Identifying such a problem in an
astronomical context, if one exists, is itself an open question.

A second direction is to confront real hardware. All results here come from a noiseless
simulator, which is the best case for quantum methods, so a natural extension is to
repeat the comparisons on physical devices together with error mitigation, both to
measure the cost of hardware noise and to test whether the structural properties seen in
simulation, such as the smoother feature maps of the quanvolutional filter, survive in
practice.

Finally, the comparisons here were bounded by eight qubits, a single scientific domain
and a handful of hard class pairs. Larger qubit counts as simulators and hardware
improve, and the same fair-comparison framework applied to other signals, would show
whether the parity found here is specific to SDSS spectra or a more general feature of
near-term quantum machine learning. The geometric-difference pre-screening used in the
kernel experiment is a cheap and principled way to decide in advance which of these
problems are even worth the effort. None of these directions is likely to overturn the
present finding on today's hardware, but each would sharpen the question of where, if
anywhere, a near-term quantum advantage on real scientific data can be found.
